\chapter{实验结果与分析}\label{chap:results}

本章以闭环成功率为主指标，报告动态抓取实验结果。除特别说明外，主表中的每个策略--环境组合均包含750个闭环评估回合。所有实验均在随机姿态翻滚平移场景中进行——立方体初始姿态均匀随机，绕空间轴以$15\,\mathrm{deg/s}$旋转并同时平移。为验证基线策略实现的正确性，在较简单的固定轴自旋平移场景中进行补充验证（详见\S\ref{sec:main-clean}）。本章的分析顺序与论文主线一致：首先比较无噪观测下的策略性能，然后分析目标中心表示，接着讨论噪声鲁棒性——包括在线观测噪声增强的收益与代价、专家FSM-MPC控制器的卡尔曼滤波评估、学习策略在显式状态估计下的表现，随后分析数据量、历史长度和模型容量边界。

\section{无噪观测下的策略比较}\label{sec:main-clean}

表~\ref{tab:main-clean-comparison}报告了四类策略的闭环成功率，图~\ref{fig:ch4-clean-main}以柱状图呈现相同数据。所有行均使用完整目标中心表示、历史长度16和MPC跟踪执行器。

\begin{table}[!htbp]
    \caption{主实验结果。所有结果均为750回合闭环成功率；方括号内为95\% Wilson置信区间。}
    \label{tab:main-clean-comparison}
    \centering
    \small
    \renewcommand{\arraystretch}{1.12}
    \begin{tabularx}{\textwidth}{>{\raggedright\arraybackslash}p{0.14\textwidth}ZZZZ}
        \hline
        策略 & BC & BC-RNN & ACT & DP 12M \\
        \hline
        闭环成功率 & \ResultCI{0.1\%}{[0.0, 0.8]} & \ResultCI{48.7\%}{[45.1, 52.2]} & \BestResultCI{59.2\%}{[55.6, 62.7]} & \ResultCI{2.5\%}{[1.6, 3.9]} \\
        \hline
    \end{tabularx}
\end{table}

\begin{figure}[!htbp]
    \centering
    \includegraphics[width=0.85\textwidth]{fig_ch4_clean_main}
    \caption{主实验结果可视化。柱高为750回合闭环成功率，误差线为95\% Wilson置信区间。}
    \label{fig:ch4-clean-main}
\end{figure}

四类策略中，ACT的闭环成功率最高（59.2\%），其95\%置信区间[55.6, 62.7]与BC-RNN的[45.1, 52.2]不重叠，差异显著。BC-RNN的48.7\%远高于单步BC的0.1\% [0.0, 0.8]，说明GRU循环历史编码为策略提供了有效的时序上下文。BC几乎完全失败（750回合仅成功1次），可见单步前馈回归无法处理随机姿态翻滚平移所需的抓取面选择、姿态对齐和闭合时机判断。DP 12M仅2.5\% [1.6, 3.9]，在本文低维目标中心动作块接口、10步DDIM推理和固定噪声调度设置下未能将生成式建模能力转化为闭环成功率——此结论的范围限定参见\S\ref{sec:limitations}限制六；完整容量扩展和表示消融结果见附录\S\ref{sec:appendix-dp}。

为排除BC实现错误的可能性，本文以较简单的固定轴自旋平移场景（立方体初始姿态固定，绕$z$轴以$30$--$60\,\mathrm{deg/s}$自旋）补充评估了同一BC模型：其闭环成功率达45.7\% [42.2, 49.3]，确认BC实现正确，在随机姿态翻滚平移场景中的失效源自场景的内在难度而非代码错误。随机姿态翻滚平移的难点并不在于自旋速度（其角速度$15\,\mathrm{deg/s}$低于固定轴场景的$30$--$60\,\mathrm{deg/s}$），而在于初始姿态均匀随机和空间轴旋转使抓取面选择和姿态对齐变得复杂。

为排除DP实现错误的可能性，同一DP 12M模型在固定轴自旋平移场景中达41.9\% [38.4, 45.4]（附录表~\ref{tab:appendix-dp-scaling}），确认DP实现正确，在随机姿态翻滚平移场景中的低成功率同样源自场景的内在难度。

对失败回合的观察进一步印证了上述分析。从策略的失败评估轨迹中可归纳出三类典型失效模式。第一类是夹爪过早闭合：策略在末端执行器尚未充分到达抓取位置或姿态尚未对齐时即触发闭合，夹爪闭空后物体留在原处或被轻微碰偏——BC几乎全部失败属于此类，ACT和BC-RNN的剩余失败也以此为主。第二类是对齐阶段未对准或边缘碰撞：策略已接近物体，但闭合前夹爪的局部姿态未精确对准物体两侧，手指撞到边缘或角点将其推开或旋转脱出。第三类是阶段转换未完成的少数情况：策略未清晰地从接近转入闭合或搬运，夹爪保持打开、末端在物体附近徘徊，或阶段转换迟疑导致任务超时。

这三类失效模式从性质上印证了本文的两个核心设计动机：（a）闭合时机判断需要时序上下文——策略必须从历史观测中推断当前是否已对齐足够久、可以安全闭合；（b）相对姿态的精确对齐需要通过目标中心坐标来组织多步动作——这正是目标中心动作块接口设计的初衷。

\section{目标中心表示消融}

表~\ref{tab:representation-ablation}比较目标中心表示和世界坐标表示，图~\ref{fig:ch4-representation-ablation}以热力图展示表示差异的整体结构。本节直接检验本文方法的核心设计：策略是否应学习相对目标的抓取动作，而不是世界坐标下的绝对末端参考。

\begin{table}[!htbp]
    \caption{表示消融结果。所有结果均为750回合闭环成功率；方括号内为95\% Wilson置信区间。}
    \label{tab:representation-ablation}
    \centering
    \small
    \begin{tabular}{lc}
        \hline
        表示 & ACT \\
        \hline
        紧凑目标中心 & \ResultCI{57.3\%}{[53.8, 60.8]} \\
        完整目标中心 & \BestResultCI{59.2\%}{[55.6, 62.7]} \\
        紧凑世界坐标 & \ResultCI{8.8\%}{[7.0, 11.0]} \\
        完整世界坐标 & \ResultCI{27.2\%}{[24.1, 30.5]} \\
        \hline
    \end{tabular}
\end{table}

\begin{figure}[!htbp]
    \centering
    \includegraphics[width=0.85\textwidth]{fig_ch4_representation_ablation}
    \caption{表示消融结果热力图。横轴为紧凑/完整表示，纵轴为目标中心/世界坐标系；每个单元格为750回合闭环成功率（\%）。目标中心表示整体优于世界坐标表示，差距超过Wilson置信区间宽度。}
    \label{fig:ch4-representation-ablation}
\end{figure}

表示消融的结果比架构排序更为稳定：目标中心表示明显优于世界坐标表示。ACT使用目标中心表示时达到57.3\%--59.2\%，而使用世界坐标表示时仅为8.8\%--27.2\%，置信区间完全分离——表示选择造成的成功率差距（约30--50个百分点）远超过ACT与BC-RNN之间的架构差距（约10个百分点，见\S\ref{sec:main-clean}）。因此，在本文的专家--学生框架中，将全局平移和姿态变化吸收进相对表示中是比模型选择更为根本的设计决策。

\section{噪声鲁棒性与状态估计}

\subsection{在线观测噪声增强训练}

图~\ref{fig:ch4-noise-robustness}展示无噪声训练策略和在线观测噪声增强策略在噪声观测下的表现。在线观测噪声增强在训练时注入与评估相近的观测扰动，以检验噪声增强能否提升闭环鲁棒性。

\begin{figure}[!htbp]
    \centering
    \includegraphics[width=0.85\textwidth]{fig_ch4_noise_robustness}
    \caption{噪声观测下的闭环成功率变化。每个策略内的两根柱分别表示无噪声训练版本和在线观测噪声增强训练版本；柱高为750回合成功率，误差线为95\% Wilson置信区间，标注为噪声增强相对无噪声训练的点估计差值。}
    \label{fig:ch4-noise-robustness}
\end{figure}

噪声观测下所有策略的成功率均大幅低于无噪观测。无噪声训练的BC-RNN仅达到16.5\%，ACT仅2.1\%。在线观测噪声增强训练的效果因策略而异：ACT从2.1\%提升至26.0\%（置信区间不重叠，提升显著）；BC-RNN从16.5\%变为18.1\%（置信区间重叠，变化不显著）。

然而，噪声增强训练对无噪性能存在代价：ACT从无噪声训练的59.2\% [55.6, 62.7]降至噪声增强训练的45.5\% [41.9, 49.1]（置信区间不重叠，下降显著）；BC-RNN也从48.7\% [45.1, 52.2]降至34.9\% [31.6, 38.4]（置信区间不重叠）。因此，噪声增强训练本质上是策略相关的鲁棒性权衡：ACT噪声观测下提升显著（$+23.9$个百分点），但无噪性能下降同样显著（$-13.7$个百分点）；BC-RNN的噪声观测收益仅为$+1.6$个百分点且置信区间重叠，同时无噪性能也下降了$-13.8$个百分点。

\subsection{显式状态估计与隐式时间冗余利用的互补作用}

噪声增强训练的有限收益指向第一章提出的核心问题：动态抓取中的噪声鲁棒性应来自显式状态估计（如卡尔曼滤波）还是学习策略自身的隐式时间冗余利用能力。本节从两个方向检验状态估计的作用：首先评估专家FSM-MPC控制器在无滤波和有卡尔曼滤波下的强噪声表现，然后评估学习策略在卡尔曼滤波观测下的闭环成功率变化。

专家FSM-MPC控制器接收与学习策略相同的噪声感知流（漂移、高斯抖动、离群点、丢帧和时延叠加），不经过任何学习策略。无滤波时，专家在强噪声下完全失败（0/750，95\% CI [0.0, 0.5]）。可见基于手工规则和MPC的专家控制器对感知噪声极其敏感：阶段判定、抓取面选择和参考轨迹构造都会因观测扰动而失效。

当在专家控制器前端加入卡尔曼滤波器后，闭环成功率从0.0\%恢复到56.4\%（95\% CI [52.8, 59.9]）。这一巨大恢复（56.4个百分点）说明：依赖手工阶段规则和MPC优化来构造参考轨迹的控制器，显式状态估计是必要的；没有滤波，感知噪声会系统性地破坏阶段逻辑和跟踪参考。

然而，学习策略是否同样依赖显式状态估计？图~\ref{fig:ch4-kalman-state-estimation}将专家控制器、BC-RNN和ACT在滤波前后的成功率放入同一张配对柱状图比较：在学习策略前端加入卡尔曼滤波器，检验显式滤波对已具备历史窗口和动作块结构的时间策略是否仍有增益。实验使用在线观测噪声增强训练的策略检查点；直接噪声观测评估和卡尔曼滤波评估均为750回合。

\begin{figure}[!htbp]
    \centering
    \includegraphics[width=\textwidth]{fig_ch4_kalman_state_estimation}
    \caption{显式状态估计对专家控制器和学习策略的不同作用。图中三组分别为FSM-MPC专家、BC-RNN和ACT；每组左柱为原始/直接噪声观测，右柱为卡尔曼滤波观测。柱高为闭环成功率，误差线为95\% Wilson置信区间，组内标注为卡尔曼滤波观测相对原始/直接噪声观测的点估计变化。FSM-MPC两柱、BC-RNN两柱和ACT两柱均为750回合评估。}
    \label{fig:ch4-kalman-state-estimation}
\end{figure}

两组实验结果形成清晰的互补关系。FSM-MPC专家加入显式滤波后获得56.4个百分点的巨大恢复，证实了基于模型的控制器对状态估计的高度依赖。相比之下，学习策略的变化幅度远小于专家：BC-RNN从18.1\%降至12.9\%（$-5.2$个百分点），ACT从26.0\%升至28.5\%（$+2.5$个百分点）。两组学习策略的卡尔曼滤波置信区间与直接噪声观测置信区间均存在不同程度的重叠，表明卡尔曼滤波对学习策略的效应不显著或仅具有边际趋势，而非与专家恢复同等级的确定性增益。

对比结果表明：学习策略借助历史窗口和动作块结构取得了隐式时间冗余利用能力，使得显式状态估计的边际收益较小且具有策略依赖性；滤波既可能为ACT提供有限帮助，也可能因相位滞后或分布偏移降低BC-RNN表现。动态抓取中的噪声鲁棒性因此分为两条互补路径：模型控制器依赖显式状态估计，而学习策略可以利用隐式时间冗余。这一结论并非否定卡尔曼滤波的价值，而是说明状态估计策略应与控制器的结构特征匹配。

\section{数据量、历史长度和模型容量的边界}

表~\ref{tab:data-history-capacity}汇总了数据量、历史长度相关的边界实验结果，图~\ref{fig:ch4-data-history-bounds}展示数据量和历史长度趋势。数据量显著提升了ACT的性能：从10\%的44.0\% [40.5, 47.6]提升至100\%的59.2\% [55.6, 62.7]（置信区间不重叠）；对BC-RNN的提升同样显著：从10\%的26.9\%提升至100\%的48.7\%（置信区间不重叠）。

历史长度的影响并非单调。紧凑目标中心表示中，BC-RNN从历史长度1的53.9\% [50.3, 57.4]显著退化至历史长度64的38.4\% [35.0, 41.9]（置信区间不重叠），这说明过长的历史窗口可能引入不相关的远期观测，干扰GRU隐藏状态的表示质量。ACT在历史长度8达到57.5\% [53.9, 61.0]，与历史长度1的50.8\% [47.2, 54.4]的置信区间存在部分重叠，其峰值效应的解释需慎重。

\begin{table}[!htbp]
    \caption{数据量、历史长度和容量边界的代表性结果。方括号内为95\% Wilson置信区间。}
    \label{tab:data-history-capacity}
    \centering
    \small
    \begin{tabular}{lccc}
        \hline
        实验组 & 设置 & BC-RNN & ACT \\
        \hline
        数据量 & 10\% & \ResultCI{26.9\%}{[23.9, 30.2]} & \ResultCI{44.0\%}{[40.5, 47.6]} \\
        & 50\% & \ResultCI{46.9\%}{[43.4, 50.5]} & \ResultCI{52.8\%}{[49.2, 56.3]} \\
        & 100\% & \ResultCI{48.7\%}{[45.1, 52.2]} & \ResultCI{59.2\%}{[55.6, 62.7]} \\
        \hline
        历史长度 $H_o$ & 1 & \ResultCI{53.9\%}{[50.3, 57.4]} & \ResultCI{50.8\%}{[47.2, 54.4]} \\
        & 8 & \ResultCI{51.1\%}{[47.5, 54.6]} & \BestResultCI{57.5\%}{[53.9, 61.0]} \\
        & 64 & \ResultCI{38.4\%}{[35.0, 41.9]} & \ResultCI{52.7\%}{[49.1, 56.2]} \\
        \hline
    \end{tabular}
\end{table}

\begin{figure}[!htbp]
    \centering
    \includegraphics[width=\textwidth]{fig_ch4_data_history_bounds}
    \caption{数据量和历史长度边界。（a）不同训练数据比例下BC-RNN和ACT的闭环成功率。（b）紧凑目标中心表示下，不同历史长度对BC-RNN和ACT的影响。误差线为95\% Wilson置信区间。}
    \label{fig:ch4-data-history-bounds}
\end{figure}

综合来看，在ACT和BC-RNN已取得有效闭环控制的条件下，单纯增加数据量或历史长度等单一资源并不能持续带来收益。更多数据对ACT和BC-RNN均有显著帮助，但更长历史并不总能提升性能（BC-RNN在历史长度64时显著退化）。因此，制约动态抓取成功率的并非数据规模，而是表示、时间结构和执行机制的设计选择。

\section{本章小结}

本章结果可归纳出三个主要结论。第一，目标中心表示是本实验中最稳定的设计收益来源：在随机姿态翻滚平移场景中，目标中心与世界坐标之间的成功率差距（对ACT约为30--50个百分点）远超过ACT与BC-RNN之间的架构差距（约10个百分点），表示的效应在统计上显著且幅度超过架构差异。第二，随机姿态翻滚平移场景需要时间结构和阶段约束：单步BC几乎完全失败（0.1\%），BC-RNN和ACT的成功率与之置信区间完全分离（48.7\%和59.2\%）；ACT在仅使用10\%训练数据时仍达到44.0\%，说明少量数据配合时间动作块结构和变换器条件建模即可获得有效的抓取策略。第三，噪声鲁棒性分为两条互补路径：基于模型的FSM-MPC控制器严重依赖显式卡尔曼滤波（恢复56.4个百分点），而学习策略借助历史窗口和动作块结构取得了隐式时间冗余利用能力，显式滤波对学习策略的影响较小且依赖策略（BC-RNN为$-5.2$个百分点，ACT为$+2.5$个百分点）；在线观测噪声增强训练的效果同样因策略而异——ACT噪声观测下从2.1\%显著提升至26.0\%，但无噪性能下降显著，BC-RNN的噪声收益不显著且无噪性能下降。由此，在FSM-MPC专家示教框架内，动态目标抓取的闭环成功率是表示、时间动作块、阶段结构和状态估计策略共同作用的结果。
